% =====================================================================
%  Liga os metadados do latexgen.config.json aos comandos da abntex2.
%
%  As macros \lgAlgumaCoisa vem de config/metadata.tex, que e gerado.
%  \lgIfAlgumaCoisa{ha valor}{esta vazio} permite omitir blocos inteiros
%  quando o campo nao foi preenchido.
%
%  Normalmente nao ha motivo para editar este arquivo: para mudar um
%  valor, edite latexgen.config.json; para mudar a aparencia, use
%  config/style.tex.
% =====================================================================

\titulo{\lgTitle}
% A NBR 6022 pede o titulo tambem em lingua estrangeira. No modo article a
% abntex2 sempre le \theforeigntitle no \maketitle, entao este comando e
% obrigatorio mesmo quando o campo esta vazio.
\tituloestrangeiro{\lgTitleEn}
\autor{\lgAuthor}
\local{\lgCity}
\data{\lgYear}
\instituicao{\lgInstitution}
\tipotrabalho{Artigo}

% Filiacao e orientacao, impressas logo abaixo do autor.
% \maketitlehookc e o gancho da memoir executado entre o autor e a data;
% no modo article a abntex2 ja usa o gancho b para o titulo estrangeiro.
% Cada linha so aparece se o campo correspondente estiver preenchido.
\renewcommand{\maketitlehookc}{%
  \lgIfInstitution{\begin{center}\normalsize\lgInstitution\end{center}}{}%
  \lgIfAdvisor{\begin{center}\normalsize Orientador: \lgAdvisor\end{center}}{}%
}

% Metadados do PDF, usados por leitores e buscadores.
\makeatletter
\hypersetup{
  pdftitle={\@title},
  pdfauthor={\@author},
  pdfsubject={\lgKeywords},
  pdfkeywords={\lgKeywords},
  pdfcreator={latexgen}
}
\makeatother
